\def\tituloconcepto{Pulsos optimizados con GRAPE}
% Si el capitulo se compila solo (sin main.tex), esta cabecera abre el documento.
% Antes de % Si el capitulo se compila solo (sin main.tex), esta cabecera abre el documento.
% Antes de % Si el capitulo se compila solo (sin main.tex), esta cabecera abre el documento.
% Antes de \input{cabecera} cada capitulo define \tituloconcepto.
\ifdefined\enmain\else
  \documentclass[11pt,a4paper]{article}
  \input{preambulo}
  \begin{document}
  \title{\tituloconcepto}
  \author{\autor}
  \date{\today}
  \maketitle
\fi
 cada capitulo define \tituloconcepto.
\ifdefined\enmain\else
  \documentclass[11pt,a4paper]{article}
  % Preambulo compartido por todos los documentos. Compilar con XeLaTeX/LuaLaTeX
% (o Tectonic). Si se usa Overleaf, elegir el compilador XeLaTeX.
\usepackage[spanish,es-noshorthands]{babel}
\usepackage{amsmath,amssymb}
\usepackage{booktabs,array}
\usepackage[a4paper,margin=2.5cm]{geometry}
\usepackage{xcolor}
\usepackage{listings}
\usepackage[most]{tcolorbox}
\usepackage{quantikz}
\usepackage[hidelinks]{hyperref}

% ---- Autor (editar a gusto) -------------------------------------------------
\newcommand{\autor}{Genesis Silesky}

% ---- Notacion de Dirac y atajos ---------------------------------------------
% quantikz ya define \ket y \bra; DeclareRobustCommand los redefine sin error.
\DeclareRobustCommand{\ket}[1]{\left|#1\right\rangle}
\DeclareRobustCommand{\bra}[1]{\left\langle#1\right|}
\DeclareRobustCommand{\braket}[2]{\left\langle#1\middle|#2\right\rangle}
\newcommand{\C}{\mathbb{C}}
\newcommand{\R}{\mathbb{R}}
\DeclareMathOperator{\Tr}{Tr}

% ---- Cajas de apoyo --------------------------------------------------------
% "concepto": definiciones previas que hacen falta para entender el capitulo.
\newtcolorbox{concepto}[1]{breakable,colback=blue!4,colframe=blue!55!black,
  fonttitle=\bfseries,title={Concepto previo: #1}}
% "nota": aclaraciones y advertencias.
\newtcolorbox{nota}[1]{breakable,colback=orange!6,colframe=orange!70!black,
  fonttitle=\bfseries,title={#1}}

% ---- Estilos de listados ----------------------------------------------------
\lstdefinestyle{py}{language=Python,basicstyle=\ttfamily\footnotesize,
  breaklines=true,postbreak=\mbox{\textcolor{gray}{$\hookrightarrow$}\space},
  numbers=left,numberstyle=\tiny\color{gray},frame=single,
  commentstyle=\color{green!40!black}\itshape,
  keywordstyle=\color{blue!70!black}\bfseries,
  stringstyle=\color{red!60!black},showstringspaces=false,tabsize=4,
  columns=fullflexible,keepspaces=true}
\lstdefinestyle{salida}{basicstyle=\ttfamily\footnotesize,breaklines=true,
  frame=leftline,backgroundcolor=\color{gray!8},columns=fullflexible,
  keepspaces=true}

  \begin{document}
  \title{\tituloconcepto}
  \author{\autor}
  \date{\today}
  \maketitle
\fi
 cada capitulo define \tituloconcepto.
\ifdefined\enmain\else
  \documentclass[11pt,a4paper]{article}
  % Preambulo compartido por todos los documentos. Compilar con XeLaTeX/LuaLaTeX
% (o Tectonic). Si se usa Overleaf, elegir el compilador XeLaTeX.
\usepackage[spanish,es-noshorthands]{babel}
\usepackage{amsmath,amssymb}
\usepackage{booktabs,array}
\usepackage[a4paper,margin=2.5cm]{geometry}
\usepackage{xcolor}
\usepackage{listings}
\usepackage[most]{tcolorbox}
\usepackage{quantikz}
\usepackage[hidelinks]{hyperref}

% ---- Autor (editar a gusto) -------------------------------------------------
\newcommand{\autor}{Genesis Silesky}

% ---- Notacion de Dirac y atajos ---------------------------------------------
% quantikz ya define \ket y \bra; DeclareRobustCommand los redefine sin error.
\DeclareRobustCommand{\ket}[1]{\left|#1\right\rangle}
\DeclareRobustCommand{\bra}[1]{\left\langle#1\right|}
\DeclareRobustCommand{\braket}[2]{\left\langle#1\middle|#2\right\rangle}
\newcommand{\C}{\mathbb{C}}
\newcommand{\R}{\mathbb{R}}
\DeclareMathOperator{\Tr}{Tr}

% ---- Cajas de apoyo --------------------------------------------------------
% "concepto": definiciones previas que hacen falta para entender el capitulo.
\newtcolorbox{concepto}[1]{breakable,colback=blue!4,colframe=blue!55!black,
  fonttitle=\bfseries,title={Concepto previo: #1}}
% "nota": aclaraciones y advertencias.
\newtcolorbox{nota}[1]{breakable,colback=orange!6,colframe=orange!70!black,
  fonttitle=\bfseries,title={#1}}

% ---- Estilos de listados ----------------------------------------------------
\lstdefinestyle{py}{language=Python,basicstyle=\ttfamily\footnotesize,
  breaklines=true,postbreak=\mbox{\textcolor{gray}{$\hookrightarrow$}\space},
  numbers=left,numberstyle=\tiny\color{gray},frame=single,
  commentstyle=\color{green!40!black}\itshape,
  keywordstyle=\color{blue!70!black}\bfseries,
  stringstyle=\color{red!60!black},showstringspaces=false,tabsize=4,
  columns=fullflexible,keepspaces=true}
\lstdefinestyle{salida}{basicstyle=\ttfamily\footnotesize,breaklines=true,
  frame=leftline,backgroundcolor=\color{gray!8},columns=fullflexible,
  keepspaces=true}

  \begin{document}
  \title{\tituloconcepto}
  \author{\autor}
  \date{\today}
  \maketitle
\fi


\section{Pulsos optimizados con GRAPE: control como problema de optimización}

\subsection{Objetivo}
Diseñar con un algoritmo numérico (GRAPE) la forma de un pulso de
radiofrecuencia que prepare un estado de Bell en dos espines acoplados, y
comprobar cómo se comporta frente a errores de frecuencia y de amplitud.

\subsection{Conceptos previos}
\begin{concepto}{control con Hamiltoniano, fidelidad y gradiente}
Ver capítulo~12 (control con pulsos y acoplamiento), capítulo~13 (fidelidad) y
capítulo~19 (optimización con gradiente). La \emph{fidelidad} con un estado
objetivo puro $\ket\phi$ es $F=|\langle\phi|\psi\rangle|^2$ (vale $1$ si son
iguales).
\end{concepto}

\subsection{Fundamento teórico}
Hasta ahora los pulsos tenían forma rectangular (\emph{pulsos duros}: amplitud
constante durante cierto tiempo). Pero la forma de un pulso se puede elegir.
Según las notas del curso, los \emph{pulsos con forma} (gaussianos, de Hermite)
reducen los errores de excitación fuera de banda, y los algoritmos de
optimización numérica como \textbf{GRAPE} y \textbf{SMP} ajustan la forma de los
pulsos de control para maximizar la fidelidad de una operación dada: tratan el
control como un problema de optimización, no de diseño manual. La guía de la
práctica pide preparar un estado de Bell con GRAPE y hacer el pulso resistente a
errores de $\pm10$~Hz en la frecuencia y $\pm5$\,\% en la amplitud.

\paragraph{El problema.} Dos espines con acoplamiento $J=696$~Hz y control de
radiofrecuencia en $x$ e $y$ sobre cada uno:
\begin{equation}
  H(t)=\frac{2\pi J}{4}\sigma_z\otimes\sigma_z
  +\frac{2\pi}{2}\sum_{j=1}^{4}u_j(t)\,\sigma_j ,
\end{equation}
con cuatro controles $u_j(t)$ (en Hz): $\sigma_x$ y $\sigma_y$ de cada espín. Se
busca la forma de $u_j(t)$ que lleve $\ket{00}$ a $\ket{\Phi^+}=(\ket{00}+\ket{11})/\sqrt2$.

\paragraph{GRAPE.} (\emph{Gradient ascent pulse engineering.}) La duración total
$T$ se divide en $N$ pasos de amplitud constante, y la amplitud de cada control
en cada paso es un parámetro. Se maximiza
\begin{equation}
  F=\big|\langle\Phi^+|\,U_N\cdots U_2U_1\ket{00}\big|^2,\qquad
  U_k=e^{-iH_k\,\Delta t}.
\end{equation}
El gradiente respecto de cada amplitud se obtiene propagando el estado hacia
adelante y el objetivo hacia atrás (en primera aproximación en $\Delta t$), y un
optimizador (aquí L-BFGS-B) lo usa para mejorar los parámetros.

\paragraph{Robustez.} Para que el pulso resista errores se optimiza el
\emph{promedio} de la fidelidad sobre varias condiciones de error a la vez:
desintonías $\pm\Delta$ en ambos espines y factores de escala de amplitud.

\begin{nota}{Fuentes y alcance}
Este capítulo se basa en el resumen de las notas (qué son GRAPE y SMP) y en la
guía de la práctica. Los archivos de teoría de pulsos con forma y de
optimización numérica son escaneos sin texto legible. Por eso aquí \textbf{solo se
implementa GRAPE} con el estado de Bell; no se implementan SMP, los pulsos duros
de comparación ni los pulsos gaussianos o de Hermite.
\end{nota}

\subsection{Código (completamente comentado)}
\lstinputlisting[style=py]{../../fuentes/23_pulsos_grape.py}

\subsection{Explicación celda por celda}
\begin{description}
  \item[Celda 2] Define $H$ con deriva, controles, $T=1$~ms, $N=20$ pasos,
    amplitud máxima de $5$~kHz y el estado objetivo.
  \item[Celda 3] GRAPE: propagadores, fidelidad, gradiente y la función
    \texttt{optimizar}, con cuatro intentos desde puntos iniciales distintos.
  \item[Celda 4] Optimiza el pulso del caso ideal (sin errores).
  \item[Celda 5] Evalúa ese pulso con desintonías de $\pm10$~Hz y errores de
    amplitud de $\pm5$\,\% (nueve combinaciones).
  \item[Celda 6] Optimiza el promedio sobre esas nueve combinaciones.
  \item[Celda 7] Repite con errores mayores ($\pm300$~Hz y $\pm20$\,\%).
\end{description}

\subsection{Salida obtenida}
\lstinputlisting[style=salida]{../../fuentes/salidas/23_pulsos_grape.txt}
Es determinista (semillas fijas).

\subsection{Interpretación}
El pulso GRAPE prepara el estado de Bell con fidelidad $0.9979$, usando como
máximo $1694$~Hz de amplitud. Con los errores de la guía ($\pm10$~Hz y
$\pm5$\,\%), el pulso sin robustez ya mantiene una fidelidad promedio de
$0.9939$ (mínima $0.9889$), y el robusto $0.9927$ (mínima $0.9902$):
\textbf{casi no hay diferencia}, porque el pulso dura solo $1$~ms y esos errores
son pequeños. Con errores mayores ($\pm300$~Hz y $\pm20$\,\%) la diferencia es
grande: la fidelidad promedio del pulso sin robustez cae a $0.4130$ (mínima
$0.0809$), mientras que el pulso optimizado para esos errores mantiene $0.9520$
(mínima $0.8751$). Es el costo típico: el pulso robusto es un poco menos preciso
en el caso ideal ($0.9966$ frente a $0.9979$).

\subsection{Ejercicios}
\begin{enumerate}
  \item Aumenta $T$ a $2$~ms y mira cómo cambia la fidelidad y la robustez.
  \item Reduce la amplitud máxima a $1$~kHz: ¿se sigue alcanzando el Bell?
  \item Cambia el estado objetivo a $\ket{\Psi^+}=(\ket{01}+\ket{10})/\sqrt2$.
\end{enumerate}

\subsection{Errores comunes}
\begin{itemize}
  \item Creer que el optimizador siempre halla el máximo global: GRAPE puede
    quedarse en un máximo local, por eso se prueban varios puntos iniciales.
  \item Olvidar que la fidelidad reportada se calcula con el propagador exacto,
    mientras que el gradiente usa una aproximación de primer orden.
\end{itemize}

\ifdefined\enmain\else
  \end{document}
\fi

